\unnumbered{Аннотация}
Исследуется подъём маятника на тележке из нижней области и его последующее удержание сверху при ограничениях движения. В идеализированной модели с управляемым ускорением тележки сравниваются четыре алгоритма обучения с подкреплением и классический траекторный регулятор. Защитный фильтр проверяет запрос контроллера и при необходимости меняет ускорение, чтобы сохранить тормозной запас тележки. Успех удержания и соблюдение ограничений оцениваются отдельно.

Исследованы 33 обучения по 100\,000 переходов: основная серия, ограниченный подбор настроек и их проверка на новых реализациях случайности. Два из изученных методов с фильтром решают все двадцать проверочных сценариев в трёх исходных обучениях. Увеличенный начальный сбор опыта улучшает результат отбора двух неуспешных процедур, но не проходит заранее установленный критерий подтверждения на новых обучениях. Показан случай роста суммарной награды при ухудшении удержания из-за скорости тележки.

Отдельно семь закреплённых контроллеров испытаны на 140 общих стартовых состояниях: 1960 эпизодов, 980 пар с фильтром и без него. Число успехов увеличивается с 441/980 до 796/980; наблюдённые превышения желаемой границы с допуском исчезают. Однако 184/980 эпизода с фильтром не выполняют заключительное удержание, хотя достигают горизонта. Выводы относятся к фиксированному диагностическому набору и симуляции. Независимая итоговая оценка и аппаратные испытания не проводились.

\medskip\noindent\textbf{Ключевые слова:} перевёрнутый маятник; обучение с подкреплением; управление с ограничениями; защитный фильтр; повторяемость обучения; начальные условия; вычислительный эксперимент.

\unnumbered{Введение}
Перевёрнутый маятник на тележке представляет собой простую по устройству, но содержательную задачу управления. Привод перемещает тележку, а маятник вращается свободно: отдельного управляющего воздействия на его оси нет. Для подъёма из нижнего положения нужно передать маятнику энергию движением шарнира. Вблизи верхнего неустойчивого равновесия задача меняется: необходимо погасить угловую скорость и компенсировать отклонения. Ограниченная длина направляющей связывает эти этапы. Движение, полезное для подъёма, может не оставить тележке места для торможения.

При известной модели эту задачу можно решать с помощью планирования траектории и обратной связи. Альтернативу даёт обучение с подкреплением: управляющая программа улучшает поведение по опыту взаимодействия со средой. Существование обоих подходов делает маятник удобным объектом сравнительного исследования. Практический вопрос состоит не только в том, способен ли алгоритм однажды поднять маятник, но и в том, повторяется ли результат обучения, сохраняется ли он при других стартах и соблюдаются ли ограничения движения.

\textbf{Цель работы} — исследовать качество подъёма и удержания перевёрнутого маятника классическим регулятором и выбранными методами обучения с подкреплением в общей модели с ограничениями. Объект исследования — система «тележка — маятник» с управляемым ускорением тележки. Предмет исследования — результативность управления, повторяемость обучения и чувствительность закреплённых контроллеров к начальным условиям и защитному фильтру.

Для достижения цели решаются следующие задачи:
\begin{enumerate}
\item согласовать модель, ограничения, наблюдение агента и физический критерий успешного удержания;
\item организовать сравнение DQN, DDPG, SAC, TQC и классического траекторного регулятора при явно заданных условиях;
\item проверить ограниченный подбор настроек и повторяемость выбранного улучшения на новых запусках обучения;
\item оценить одни и те же контроллеры на структурированном наборе начальных состояний с фильтром и без него;
\item разделить причины неуспеха по наблюдаемым признакам и установить границы выводов.
\end{enumerate}

В обучении с подкреплением (reinforcement learning, RL) \emph{агент} получает наблюдение и выбирает действие. Правило такого выбора называется \emph{политикой}. Среда исполняет действие и возвращает новое наблюдение и численную \emph{награду}. Одна попытка управления до заданного времени или досрочной остановки называется \emph{эпизодом}. Агент обучается увеличивать накопленную награду, однако её рост не обязательно означает выполнение физической задачи. Поэтому награда, успех удержания и соблюдение ограничений в работе измеряются раздельно.

Случайная инициализация нейросетей и сбор опыта влияют на найденную политику. Начальное значение генераторов случайных чисел (\emph{seed}) задаёт конкретную реализацию случайности. Разные seed обучения не означают разные границы распределения начальных состояний. Сохранённые параметры модели образуют \emph{контрольную точку} (checkpoint). В работе различаются лучшая по правилу отбора точка и последние веса одного обучения; это две версии одной модели, а не два независимых запуска.

Для выбора модели используется \emph{валидация}: оценка на фиксированных состояниях, результаты которой участвуют в отборе. Такие состояния нельзя затем считать независимым итоговым тестом. Аналогично, проверка выбранной настройки на новых seed характеризует повторяемость обучения, но не заменяет испытание на новых физических стартах. Под \emph{робастностью} далее понимается эмпирическое сохранение качества при изменении начальных условий. Устойчивость к шумам, задержкам и ошибкам параметров в этот эксперимент не входит.

Защитный фильтр проверяет запрос контроллера и при необходимости меняет ускорение, сохраняя тормозной запас тележки. При включённом фильтре оценивается система «контроллер вместе с фильтром»: её успех нельзя целиком приписывать нейросети. В то же время защита положения не обеспечивает подъём и удержание сама по себе. Именно поэтому сравнение проводится как между процедурами обучения, так и между двумя режимами исполнения одних весов.

Эксперименты включают основную серию из 15 обучений, этап подбора настроек из 12 обучений и проверку выбранных настроек из шести обучений. Каждый запуск содержит 100\,000 переходов, суммарно — 3,3 млн. Дополнительно семь фиксированных контроллеров оценены на 140 общих стартовых состояниях: 1960 эпизодов, образующих 980 пар с фильтром и без него. Повторное использование стартов обеспечивает сопоставимость, но не превращает эти пары в независимую выборку из всего пространства состояний.

\textbf{Основные результаты.} SAC и TQC с фильтром решают все двадцать валидационных сценариев в каждом из трёх исходных обучений. Подбор увеличенного начального сбора опыта для DDPG с фильтром и SAC без него улучшает результат отбора, но не проходит последующую проверку повторяемости. В структурированной матрице включение фильтра увеличивает число успехов с 441/980 до 796/980 и устраняет наблюдённые превышения желаемой границы с учётом допуска. При этом 184/980 эпизода с фильтром завершают горизонт без успешного заключительного удержания. Эти результаты показывают различие между допустимым движением тележки и выполнением полной задачи.

\textbf{Вклад проекта} состоит в построении согласованной экспериментальной цепочки на основе лабораторного симулятора и в сравнительном анализе полученных результатов. Подготовлены интерфейс среды, обработка событий ограничений, защитный слой, протокол выбора моделей и средства анализа. Разделены влияние фильтра, случайности обучения и начального состояния; сохранены отрицательные результаты. Лабораторная динамика, основа классического управления и библиотечные алгоритмы обучения имеют самостоятельное происхождение. Разработка нового универсального алгоритма управления не заявляется.

Работа ограничена симуляцией с идеально задаваемым ускорением. Обученные политики на оборудовании не испытывались, независимая итоговая оценка после всех решений разработки не проводилась. Эти ограничения учитываются при обсуждении каждого результата.

